Os conjuntos de sentenças, em certo sentido, descrevem as estruturas nas quais são
conjuntamente verdadeiras; essas estruturas são seus \emph{modelos}. Reciprocamente,
se começamos com uma estrutura ou um conjunto de estruturas, podemos estar
interessados no conjunto de sentenças das quais elas são modelos; esta é a
\emph{teoria} da estrutura ou do conjunto de estruturas. Qualquer conjunto desse tipo de
sentenças tem a propriedade de que toda sentença que é consequência delas já
está no conjunto; elas são \emph{fechadas}. De modo mais geral, chamamos um
conjunto $\Gamma$ de teoria se ele é fechado sob consequência, e dizemos que
$\Gamma$ é \emph{axiomatizado} por~$\Delta$ se $\Gamma$ consiste em todas as
sentenças que são consequência de~$\Delta$.

A matemática fornece muitos exemplos de teorias, por exemplo, as teorias das
ordens lineares, dos grupos, ou teorias da aritmética, por exemplo, a teoria
axiomatizada pelos axiomas de Peano. Mas há muitos exemplos de
teorias importantes em outras disciplinas também, por exemplo, os bancos de dados
relacionais podem ser pensados como teorias, e a metafísica ocupa-se
de teorias da relação parte-todo, que podem ser axiomatizadas.

Uma questão significativa ao montar uma teoria para estudo é se
sua linguagem é expressiva o bastante para nos permitir formular tudo o que
queremos que a teoria trate, e outra é se ela é forte
o bastante para provar o que queremos que ela prove. Para \emph{expressar} uma
relação, precisamos de uma fórmula com o número necessário de variáveis
livres. Na \emph{teoria dos conjuntos}, temos apenas $\in$ como símbolo de
relação, mas ele nos permite expressar $x \subseteq y$ usando
$\lforall[u][(u \in x \lif u \in y)]$. A \emph{teoria dos conjuntos de
  Zermelo-Fraenkel}~$\Log{ZFC}$, de fato, é forte o bastante tanto para expressar
(quase) toda afirmação matemática quanto para (quase) provar todo
teorema matemático usando um punhado de axiomas e uma cadeia de
definições cada vez mais complicadas, como a de~$\subseteq$.
